\section{Introduction}\label{sec:introduction}

Uniform Lipschitz heights are among the simplest random surface models with a
hard local constraint.  Their planar fluctuations are expected to be governed
by the Gaussian free field.  We prove this for integer heights on the
triangular lattice with the two-arc boundary condition proposed by Schramm.
The conclusion concerns the field, including all joint test-function limits
and distributional tightness.

\subsection{The model and the limit}

Write
\[
 \Tri=\Z(1,0)+\Z(1/2,\sqrt3/2)
\]
for the unit triangular lattice.  Let $D\subset\R^2$ be bounded, smooth, and
simply connected, and let $a,b$ be distinct points of $\partial D$.  They
divide the boundary into open arcs $A_+$ and $A_-$.  For $\delta\downarrow0$,
let $D_\delta$ be simply connected polygonal domains made of triangles of
$\delta\Tri$, whose marked boundary parametrizations converge uniformly to
that of $(D,a,b)$.  Put the discrete marks at the midpoints of distinct
boundary edges.  The boundary vertices then split into disjoint sets
$A_{+,\delta}$ and $A_{-,\delta}$; no vertex receives two prescriptions.

Choose $h_\delta$ uniformly among the odd integer functions on the vertices
of $D_\delta$ satisfying
\begin{equation}\label{eq:odd-model}
 h_\delta=1\text{ on }A_{+,\delta},\qquad
 h_\delta=-1\text{ on }A_{-,\delta},\qquad
 |h_\delta(x)-h_\delta(y)|\in\{0,2\}\quad(x\sim y).
\end{equation}
Extend the function affinely on every triangle.  For a test function
$f\in C_c^\infty(D)$, write $h_\delta(f)=\int h_\delta(x)f(x)\,dx$ once its
support lies in $D_\delta$.  This convention is used for every smeared height
below.

Our Dirichlet Gaussian free field $\Phi_D$ has covariance
\begin{equation}\label{eq:gff-convention}
 \EE[\Phi_D(f)\Phi_D(g)]
 =\iint_{D\times D}f(x)G_D(x,y)g(y)\,dx\,dy,
 \qquad G_D=(-\Delta_D)^{-1}.
\end{equation}
In particular, the logarithmic singularity of $G_D(x,y)$ is
$(2\pi)^{-1}\log(1/|x-y|)$.  Let $\bar g$ be the bounded harmonic extension
of the boundary data $+1$ on $A_+$ and $-1$ on $A_-$.

\begin{theorem}\label{thm:main}
Let $D$ be a bounded smooth simply connected planar domain with two
distinct marked boundary points.  Let $D_\delta$ be simply connected
triangular-lattice polygonal approximations with uniformly converging
marked boundary parametrizations and edge-midpoint marks, and let
$h_\delta$ have the uniform odd-height law \eqref{eq:odd-model}.
There is a constant $\sigma>0$, depending only on the unit triangular-lattice
model, such that
\[
 h_\delta\ \Longrightarrow\ \bar g+\sigma\Phi_D
 \qquad\text{in }\mathcal D'(D).
\]
For every $f\in C_c^\infty(D)$, $\EE h_\delta(f)\to\int_D\bar g f$.
Consequently
\begin{equation}\label{eq:main-centered}
 \frac{h_\delta-\EE h_\delta}{\sigma}
 \ \Longrightarrow\ \Phi_D
 \qquad\text{in }\mathcal D'(D).
\end{equation}
All finite vectors of test-function integrals converge jointly to the
corresponding Gaussian vectors, and all their mixed moments converge.
The constant is specified by the absolutely convergent finite-volume series
in \cref{prop:normalization-series}; in the spectral convention below,
$\sigma=4\pi\sqrt{c_*}$.
\end{theorem}

It is convenient to prove the theorem for
\begin{equation}\label{eq:integer-height}
 H_\delta=(h_\delta-1)/2,
 \qquad g=0\text{ on }A_+,\quad g=-1\text{ on }A_-.
\end{equation}
Thus $H_\delta$ is integer valued and $|H_\delta(x)-H_\delta(y)|\le1$ on
edges.  Write $m$ for the bounded harmonic extension of $g$, so
$\bar g=1+2m$.  We will show that $H_\delta$ converges to
$m+\sqrt v\,\Phi_D$, where
\begin{equation}\label{eq:normalization-relations}
 v=(2\pi)^2c_*,\qquad \sigma^2=4v.
\end{equation}
There is no logarithmic rescaling of the distribution-valued field.  The
logarithmic divergence of point fluctuations is consistent with a finite
limit after smooth spatial averaging.

\subsection{History and significance}

Schramm's 2006 problem collection asks about the scaling limits of this
height function and its zero interface
\cite[Problem~2.2, p.~527]{Schramm2007}.
\Cref{thm:main} resolves the field part of that problem for smooth simply
connected domains, with the covariance convention made explicit in
\eqref{eq:gff-convention}.  The
interface assertion is a separate question.  The real-valued Lipschitz
model in Schramm's Problem~2.3, p.~528, is also distinct from the integer model
studied here.

Glazman and Manolescu proved logarithmic fluctuations for the uniform
integer Lipschitz model and developed a two-spin representation with
positive association and scale estimates for loops \cite{GM2021}.  Those
estimates are the probabilistic starting point of this paper.  Glazman and
Lammers subsequently proved delocalisation across a larger parameter range
and quantified it logarithmically.  Their 2025 article discusses the
Gaussian free field limit as a conjecture \cite[Section~1.2.4]{GL2025}.
Logarithmic fluctuations alone do not determine a Gaussian scaling limit.
The additional work here is to obtain comparisons under rare pinning,
identify the covariance, control the imposed boundary values, and determine
all joint moments.

For real-valued gradient fields with smooth uniformly convex interactions,
Miller proved a bounded-domain GFF limit with continuous perturbations of
affine boundary data \cite[Theorem~1.1]{Miller2011}.  The integer-valued law
studied here, with its hard constraint on every edge increment, lies outside
those hypotheses.

Besides the loop estimates, we use the discrete midpoint
Pr\'ekopa--Leindler inequality of Klartag--Lehec
\cite[Theorem~1.4]{KlartagLehec2019}, in the nonnegative-function form of
Gozlan--Roberto--Samson--Tetali \cite[Theorem~3]{GRST2019}.  Its role is to
turn central probability estimates into uniform higher moments.  The
rounding and parity conditions needed for this application are verified
in \cref{sec:moments}.

Reflection positivity is a classical method for converting spatial
symmetry and positivity into spectral information
\cite{FILS1978}.  The reflection quotient and positive transfer used here
are closely related to the constructions in
\cite[Sections~3.1--3.3 and 4.1]{Usui2012}.  We establish their exact
site-row and auxiliary-phase versions directly.  The final moment
argument has a related predecessor in Kenyon's identification of dimer
height moments by matching collision poles in a determinant pairing identity
\cite[Lemma~3.1 and Proposition~3.2]{Kenyon2001}; here the collision
coefficients come from local cut mixing rather than a determinant formula.

Duminil-Copin--Kozlowski--Lammers--Manolescu prove a full-plane GFF limit
for the zero-slope square-lattice six-vertex state with
$a=b=1$ and $c\in[\sqrt3,2]$
\cite[Theorems~2.2 and 2.8]{DCKLM2026}.  Their transfer representations,
analytic translation arguments, harmonic correlations and collision
analysis are substantial methodological precedents
\cite[Theorems~4.12 and 4.15; Sections~7--9]{DCKLM2026}.
The argument below uses triangular site reflections and a finite-rotation
angular sum to determine the plane covariance.  Separated rare pins and
boundary attachments then connect this information to the actual
two-arc bounded-domain law.  The normalization is specified by finite
counts and Fourier coefficients in \cref{sec:normalization}.

For comparison, the prescribed iterated balanced-torus limit of the
square-lattice six-vertex model with $a=b=1$ gives a balanced zero-slope
plane state with a GFF scaling limit for each fixed $0<c\le2$
\cite[Theorem~1.1]{OpenAIBalancedSixVertex2026}.
This plane-state result does not identify the present two-arc
triangular-lattice law or its normalization.

A companion article proves the GFF limit for weighted integer Lipschitz
heights with zero boundary, for each fixed edge weight
$x\in[1/\sqrt2,1]$ \cite[Theorem~1.2]{OpenAIWeighted2026}.
Its weighted conditional laws and moment estimates are proved there.
The present two-arc boundary theorem and normalization series are retained
separately; neither field theorem is used as a premise for the other.

\subsection{Proof strategy}

The main task is to identify every subsequential limit $H$ of the integer
height field.  A key intermediate conclusion is
\begin{equation}\label{eq:strategy-harmonicity}
 \EE\left[H(\Delta\phi)\prod_{i=1}^qH(f_i)\right]=0,
\end{equation}
for smooth compactly supported tests $\phi,f_1,\ldots,f_q$ in $D$ with
pairwise disjoint supports, including $q=0$.  This says that the moment
distributions are harmonic in each variable away from collisions.  Boundary
traces and the singularities at collisions will then determine those moments.
The difficulty is to obtain this identity for absolute heights: reflection
positivity first gives a plane identity, and moving pins initially permits
only observations of local height differences.

\paragraph{A spectral identity from reflection positivity.}
The height gradient has a stationary plane limit.  Reflection through rows
of lattice sites gives a positive normal transfer operator.  Its spectral
representation expresses the covariance as a mixture of Cauchy kernels
in the normal frequency.  The fourth angular harmonic averages to zero
under the six lattice rotations, whereas its Cauchy average is nonnegative
up to a summable lattice error.  Comparing the two gives an integrable
nonnegative dispersion defect.  This forces every small-frequency limit to have density
$c_*/|p|^2$, and a signed angular sum determines $c_*$ uniquely.
For a test $\phi$ supported on one side of a reflection line, the covariance
formula makes the covariance of the height average against $\Delta\phi$
with its reflected copy tend to zero.  This is the reflection-null identity;
it does not assert that the ordinary variance vanishes.  The signed angular
sum also gives the finite-volume normalization formula in
\cref{sec:normalization}.  The spectral argument uses positivity and the
additional rotation directions rather than an explicit solution of the model.

\paragraph{Moving rare pins and recovering the boundary.}
We encode the height modulo four by two spins taking values $\pm1$ at each vertex;
boundary pins prescribe these signs.
Boundary pins surround a typical interior insertion, so they cannot all
initially lie behind one reflection line.  We open short gaps inside
boundary windows on which only one of the two spins is pinned.
A multiplicative comparison for separated pinning
events controls the resulting reflection norms even when the pin events
are rare.  Positive transfer powers permit analytic continuation of
translations in several lattice-normal directions, adapting the method of
\cite[Section~7.1]{DCKLM2026}.  This moves the pins
past the reflection-null insertion while preserving its limiting identity
against observations of local height differences.
Circuits carrying the prescribed spin sign then close the gaps and force
the disconnected pin arcs to use one common height offset.

To recover the original boundary condition, we prove height-path estimates
at a constant boundary arc.  The key attachment estimate derives an open
level-line connection from a synthetic mean comparison, with explicit
stopping rules that preserve the Gibbs law in the unexamined region.
The estimate survives the small boundary pieces created when windows are
removed behind additional arches.  It identifies the harmonic mean and supplies deterministic
boundary reference averages.

\paragraph{Harmonic moments and Gaussianity.}
Averaging near distinct points of one constant boundary arc gives its
prescribed value in the limit, with depths sent to zero before the number
of points increases.  These reference averages fix the additive constant left
undetermined by local height differences and yield
\eqref{eq:strategy-harmonicity} for absolute observations.
The resulting moment functions have only logarithmic collision singularities.
Local cut mixing and the plane
covariance determine each singular coefficient.  Subtracting the associated
Green functions gives the Gaussian moment recursion.  Uniform moment
bounds supply both tightness and the passage from moment identities to the
centered distributional limit.

\subsection{Conventions and limit order}

The symbols $C,c,\alpha$ denote positive constants, whose values may change
between occurrences.  Dependence on a fixed finite macroscopic arrangement
is allowed when stated.  A normalized test at radius $r$ is supported in a
ball of radius $r$ and has size $O(r^{-2})$; a zero-total test has integral
zero.  The Euclidean Fourier convention is
$\widehat f(p)=\int e^{-ip\cdot x}f(x)\,dx$.

All arguments with prescribed boxes, collars, or tubes first fix their
positive widths and send the lattice mesh to zero.  Later limits that
shrink windows, refine grids, or approach the boundary are taken only
after this mesh limit.  In particular, a lower bound depending on a fixed
aspect ratio is never used as a uniform bound for a degenerating shape.
When comparison curves are random, the specified exploration is stopped
before inspecting the region whose conditional law is used.

The proof follows this order.  \Cref{sec:spin,sec:cuts,sec:moments} establish
the spin tools, conditional cuts, and moments.  \Cref{sec:spectral,sec:normalization}
identify the plane covariance and its coefficient.  \Cref{sec:pins,sec:boundary}
transfer the reflection identity and prove the boundary estimate.
\Cref{sec:limits,sec:gaussian} identify all subsequential limits and complete
the proof of \cref{thm:main}.
